% Compile twice: lualatex 20001255.tex
\documentclass[11pt]{article}
\usepackage[margin=25mm]{geometry}
\usepackage{fontspec,unicode-math}
\setmainfont{Latin Modern Roman}
\setmathfont{Latin Modern Math}
\usepackage{amsmath,amssymb,mathtools}
\usepackage{xurl,hyperref}
\hypersetup{colorlinks=true,urlcolor=blue,linkcolor=blue}
\setcounter{secnumdepth}{0}
\setlength{\parindent}{0pt}
\setlength{\parskip}{5pt}
\setlength{\emergencystretch}{3em}
\newcommand{\sref}[2]{\hyperlink{source:#1:#2}{[S#2]}}
\newcommand{\cataloguescope}[1]{\paragraph{Source question.}#1}
\newcommand{\floor}[1]{\left\lfloor#1\right\rfloor}
\begin{document}
% UnsolvedMath ID 20001255; source catalogue code AIM-COMPUTATION-0093
\setcounter{section}{20001254}
\section{Adjacent nonisogenous fibers in the Legendre family}\label{20001255}
{\small\sffamily UnsolvedMath ID 20001255 · Computer Science}
\subsection{Definitions and mathematical statement}

\paragraph{Shared notation.}Standard number systems are denoted by $\mathbb N,\mathbb Z,\mathbb Q,\mathbb R,\mathbb C$; index ranges, computational models and specialized notation follow the cited problem source and the definitions below.


An elliptic curve over $\mathbb F_p$ is a nonsingular genus-one curve with a chosen rational point. Its rational-point count is $\#E(\mathbb F_p)$. The source seeks adjacent parameters in a nonconstant family with unequal point counts, corresponding to nonisogenous fibers over the stated finite field.
\paragraph{Statement recorded in the supplied source.}
Question 5. Suppose that you have a nonconstant family of elliptic curves Eλ

over Fp (e.g. the Frey curves), λ ∈ P1(Fp). Can you find deterministically λ such that Eλ

and Eλ+1 are not isogenous (i.e. Eλ(Fp) 6 = Eλ+1 (Fp)). (Elkies) \sref{20001255}{2}
\paragraph{Scope qualification.}
The extraction drops the point-count symbols. The primary source uses point counts and remarks that p should be sufficiently large and the family degree small, without giving a numerical uniform threshold.
\paragraph{Source status.}
The supplied catalogue labels this record partially solved. Its source-status metadata was checked against the supplied snapshot on 2026-10-11. This import does not independently establish current open status or mathematical correctness.
\subsection{Short English statement}
Can one deterministically find two adjacent members of an elliptic-curve family that have different finite-field point counts?
\subsection{Sources}
\begin{itemize}
\item[\hypertarget{source:20001255:1}{S1}] Ulam AI and the UnsolvedMath Contributors, UnsolvedMath: A Curated Collection of Open Mathematics Problems, record 20001255 (AIM-COMPUTATION-0093). Catalogue curation and metadata: CC BY 4.0. The supplied source question is retained; definitions and the English summary are editorial additions.
\url{https://huggingface.co/datasets/ulamai/UnsolvedMath}.
\item[\hypertarget{source:20001255:2}{S2}] AIM Workshop Problem Lists, original problem source linked by catalogue record AIM-COMPUTATION-0093
\url{https://aimath.org/WWN/primesinp/primesinp.pdf}.
\end{itemize}
\paragraph{Preparation and review.}
Original question attributed to the cited source and UnsolvedMath contributors. Codex (GPT-6; exact serving revision not exposed) assisted with definitions, English summaries, source inspection and formatting. No new human mathematical authorship or independent proof verification is asserted. Canonical ID 20001255 and source code AIM-COMPUTATION-0093 checked against the complete supplied registry and existing repository statements; known flagged repetitions and obvious identical conjectures excluded. Full mathematical-equivalence/literature review has not been performed. Statement preparation and final review followed contribution prompts C and D; automated checks verify structure and identity, not open status.
\label{end:20001255}
\end{document}
